\begin{ouverture}[geometrie]{15}{Le cosinus dans le triangle rectangle}{Partie 4 · Géométrie et espace}
{Hypoténuse · Côté adjacent · Longueur · Angle}
\objectif{Identifier les côtés relatifs à un angle aigu.}
\objectif{Calculer une longueur à l’aide du cosinus.}
\objectif{Déterminer un angle avec la calculatrice en degrés.}
\prerequis{Pythagore, quotients et équations simples.}
\end{ouverture}
\pagestyle{otMain}\essentiel\fontsize{11}{15}\selectfont
\phantomsection\label{cours-15-1}
\essSection{1}{Un rapport associé à un angle}
\begin{definition}Dans un triangle rectangle, le cosinus d’un angle aigu est le quotient longueur du côté adjacent sur longueur de l’hypoténuse. Le côté adjacent est celui qui touche l’angle sans être l’hypoténuse.\end{definition}
\label{fig-cours-15-1-cours-cosinus}\figureGeo{cours-cosinus}

\begin{propriete}Dans ABC rectangle en A, $\cos(\widehat{ABC})=AB/BC$ et $\cos(\widehat{ACB})=AC/BC$. Pour un angle aigu, ce rapport est strictement compris entre 0 et 1. Des triangles rectangles ayant le même angle aigu donnent le même rapport, ce qui se relie à Thalès.\end{propriete}
\begin{exemple}Si $AB=6$ et $BC=10$ cm, $\cos(\widehat B)=6/10=0{,}6$. Pour l’angle C, le côté adjacent change : ce n’est plus AB mais AC. L’hypoténuse, elle, reste BC.\end{exemple}
\clearpage
\phantomsection\label{cours-15-2}
\essSection{2}{Choisir l’opération}
\begin{savoirfaire}[Calculer une longueur]Écrire d’abord le rapport avec les noms des côtés. Si l’hypoténuse $h$ et l’angle $\alpha$ sont connus, le côté adjacent vaut $h\cos\alpha$. Si l’adjacent $a$ est connu, l’hypoténuse vaut $a/\cos\alpha$.\end{savoirfaire}
\begin{exemple}Pour un angle de $40^\circ$ et une hypoténuse de 12 cm, l’adjacent vaut $12\cos40^\circ\approx9{,}19$ cm. Pour un adjacent de 7 cm au même angle, l’hypoténuse vaut $7/\cos40^\circ\approx9{,}14$ cm.\end{exemple}
\begin{savoirfaire}[Retrouver un angle]À partir de $\cos\alpha=0{,}6$, utiliser la commande cosinus inverse, souvent $\cos^{-1}$ ou arccos. En mode degrés, on trouve $\alpha\approx53{,}13^\circ$. La commande inverse ne signifie pas $1/\cos$.\end{savoirfaire}
\begin{attention}Vérifier le mode degrés de la calculatrice. La formule ne s’applique ici qu’à un triangle rectangle et à un angle aigu. Pour l’arrondi final, conserver la valeur complète du cosinus dans les calculs intermédiaires. Sinus et tangente ne sont pas nécessaires dans cette édition de Quatrième.\end{attention}

\clearpage\rubrique{Découverte}
\begin{blocLivre}{ex-15-1}\begin{exercice}[1]{Nommer les côtés}
ABC est rectangle en A. Pour l’angle B, nommer l’hypoténuse et le côté adjacent.
\label{fig-ex-15-1}\figureGeo{cos-01}
\end{exercice}\end{blocLivre}
\begin{blocLivre}{ex-15-2}\begin{exercice}[1]{Calculer un cosinus}
ABC est rectangle en A, avec $AB=6$ et $BC=10$ cm. Donner le cosinus de l’angle B.
\label{fig-ex-15-2}\figureGeo{cos-02}
\end{exercice}\end{blocLivre}
\begin{blocLivre}{ex-15-3}\begin{exercice}[1]{Changer d’angle}
ABC est rectangle en A, $AC=8$ cm et $BC=10$ cm. Donner le cosinus de l’angle C.
\label{fig-ex-15-3}\figureGeo{cos-03}
\end{exercice}\end{blocLivre}
\begin{blocLivre}{ex-15-4}\begin{exercice}[1]{Un angle de soixante degrés}
ABC est rectangle en A, $\widehat B=60^\circ$, $BC=12$ cm. Calculer AB, sachant $\cos60^\circ=0{,}5$.
\label{fig-ex-15-4}\figureGeo{cos-04}
\end{exercice}\end{blocLivre}
\begin{blocLivre}{ex-15-5}\begin{exercice}[1]{Une longueur approchée}
ABC est rectangle en A, $\widehat B=35^\circ$ et $BC=8$ cm. Calculer AB au centième.
\label{fig-ex-15-5}\figureGeo{cos-05}
\end{exercice}\end{blocLivre}

\rubrique{Entraînement}
\begin{blocLivre}{ex-15-6}\begin{exercice}[2]{Retrouver l’hypoténuse}
ABC est rectangle en A, $\widehat B=40^\circ$ et $AB=7$ cm. Calculer BC au centième.
\label{fig-ex-15-6}\figureGeo{cos-06}
\end{exercice}\end{blocLivre}
\begin{blocLivre}{ex-15-7}\begin{exercice}[2]{Déterminer un angle}
ABC est rectangle en A avec $AB=5$ et $BC=13$ cm. Calculer l’angle B au degré près.
\label{fig-ex-15-7}\figureGeo{cos-07}
\end{exercice}\end{blocLivre}
\begin{blocLivre}{ex-15-8}\begin{exercice}[2]{Un angle puis l’autre}
ABC est rectangle en A, $AB=9$ et $BC=15$ cm. Calculer les angles B et C au dixième.
\label{fig-ex-15-8}\figureGeo{cos-08}
\end{exercice}\end{blocLivre}
\begin{blocLivre}{ex-15-9}\begin{exercice}[2]{L’adjacent à C}
ABC est rectangle en A, $\widehat C=50^\circ$ et $BC=10$ cm. Calculer AC au centième.
\label{fig-ex-15-9}\figureGeo{cos-09}
\end{exercice}\end{blocLivre}
\begin{blocLivre}{ex-15-10}\begin{exercice}[2]{Une hypothèse manquante}
Un triangle ABC a $AB=6$ et $BC=10$ cm, sans information d’angle droit. Peut-on écrire $\cos\widehat B=6/10$ ?
\label{fig-ex-15-10}\figureGeo{cos-10}
\end{exercice}\end{blocLivre}

\rubrique{Problèmes}
\begin{blocLivre}{ex-15-11}\begin{exercice}[3]{Inclinaison d’une échelle}
Dans un modèle, une échelle de 4 m s’appuie sur un mur perpendiculaire au sol. Son pied est à 1,5 m du mur. Calculer son angle avec le sol au degré près et la hauteur atteinte au centimètre.
\label{fig-ex-15-11}\figureGeo{cos-11}
\end{exercice}\end{blocLivre}
\begin{blocLivre}{ex-15-12}\begin{exercice}[3]{Une rampe}
Une rampe rectiligne mesure 6 m ; sa projection horizontale mesure 5,7 m. Quel angle fait-elle avec l’horizontale, au dixième de degré ?
\label{fig-ex-15-12}\figureGeo{cos-12}
\end{exercice}\end{blocLivre}
\begin{blocLivre}{ex-15-13}\begin{exercice}[3]{Un chemin incliné}
Un chemin rectiligne de 3,2 km monte avec un angle de $27^\circ$ par rapport à l’horizontale. Calculer sa projection horizontale au centième de km.
\label{fig-ex-15-13}\figureGeo{cos-13}
\end{exercice}\end{blocLivre}
\begin{blocLivre}{ex-15-14}\begin{exercice}[3]{Comparer deux triangles}
Deux triangles rectangles ont la même hypoténuse de 10 cm. Les angles étudiés valent $30^\circ$ et $60^\circ$. Comparer les côtés adjacents à ces angles.
\label{fig-ex-15-14}\figureGeo{cos-14}
\end{exercice}\end{blocLivre}
\begin{blocLivre}{ex-15-15}\begin{exercice}[3]{Deux outils pour une hauteur}
Une tour verticale AC se trouve à 9 m du point B au sol. L’angle entre BA et BC vaut $58^\circ$. Calculer BC puis la hauteur AC, au dixième de mètre.
\label{fig-ex-15-15}\figureGeo{cos-15}
\end{exercice}\end{blocLivre}

\clearpage
